\chapter{数据工程与继续预训练：数据构建、样本拼接、增量预训练}

\begin{introduction}
\item {\heiti 本章核心问题}：为了继续改造企业知识助手，训练数据应该从哪里来、如何组织、何时用于 SFT、何时用于继续预训练？
\item {\heiti 主案例推进}：为企业知识助手建立“知识文本线”“任务样本线”和“评测回流线”三条数据资产主线。
\item {\heiti 读完后能做什么}：看到一份原始企业材料时，知道它更适合进入检索库、SFT 样本池，还是继续预训练语料池，并知道进入前要经过哪些清洗与质检步骤。
\end{introduction}

\section{学习目标}
\begin{itemize}
  \item 理解训练数据在 SFT、继续预训练和领域适配中的不同角色。
  \item 掌握样本清洗、格式化、样本拼接与增量预训练的基本策略。
  \item 能为企业知识助手建设一套可持续扩展、可追踪、可回滚的数据资产。
\end{itemize}

\section{问题背景}
模型工程的很多上限并不由训练框架决定，而由数据工程决定。没有结构清楚、边界明确的数据，微调只会把混乱放大。尤其是在企业场景中，文本来源复杂：制度文档、FAQ、工单记录、产品说明、会议纪要、客服对话都可能被拿来“喂模型”。如果不先建立数据分层，模型很快就会学到不一致风格、过期事实和模糊表达。

很多团队在这里犯的第一个错误，是把“有文本”误当成“有训练数据”。原始材料只是原料，不是天然可训练样本。第二个错误，是把所有问题都归到 SFT 或都归到继续预训练。更可靠的工程判断不是看名字，而是看症状：如果模型连领域文本都读不顺、术语分布不熟、日志和工单格式理解吃力，这更像语言分布问题；如果模型知识大体够用，但输出格式不稳、任务边界不清、口吻不符合企业要求，这更像任务行为问题。

还有一个常被忽略的事实是，所需数据量并不只由“模型多大”决定，更取决于\textbf{原模型分布与目标分布的距离}。如果企业材料和通用互联网文本风格接近，少量高质量样本就可能见效；如果材料充满内部缩略语、强结构字段、冷门术语和长尾流程，再多泛泛样本也不一定有用。

\begin{table}[H]
\centering
\small
\begin{tabularx}{\textwidth}{>{\raggedright\arraybackslash}p{3cm}>{\raggedright\arraybackslash}p{4.2cm}>{\raggedright\arraybackslash}X}
\toprule
\textbf{观察到的症状} & \textbf{更可能的问题类型} & \textbf{优先动作} \\
\midrule
回答内容基本对，但格式不稳、风格混乱 & 任务行为问题 & 先改提示模板，再补 SFT \\
制度名词、工单字段、日志格式读不顺 & 领域语言分布问题 & 优先考虑继续预训练 \\
答案依赖最新制度、最新版本说明 & 动态知识问题 & 优先进入 RAG，而不是写死到参数里 \\
多个子任务都不稳定，但都围绕同一输出规范 & 指令与样本设计问题 & 先收紧任务定义，再建设 SFT 样本 \\
\bottomrule
\end{tabularx}
\caption{SFT、继续预训练与 RAG 的症状化判断}
\label{tab:data-routing-diagnosis}
\end{table}

\section{核心概念}
\subsection{三类数据资产}
从用途上看，本书把数据分成三类：
\begin{enumerate}
  \item \textbf{继续预训练数据}：用于让模型更熟悉领域语言分布。
  \item \textbf{SFT 数据}：用于教模型完成具体任务、遵循输出格式。
  \item \textbf{评测数据}：用于检验模型和系统是否真的变好。
\end{enumerate}

这三类数据不能混为一谈。继续预训练强调规模和语料分布，SFT 强调指令与答案质量，评测数据强调代表性、稳定性和与训练链路的隔离。一本制度手册、一个 FAQ 页面、一次真实工单对话，都可能同时服务于这三类资产，但进入不同资产后的加工方式完全不同。

\subsection{先分流，再构造：原始材料不是天然训练样本}
同一份原始材料进入系统前，至少要先回答三个问题：它是用来\textbf{被检索}、\textbf{教任务行为}，还是\textbf{让模型适应语言分布}？如果这个问题没有先回答，后面的清洗、切分、训练和评测都会变得混乱。

对企业知识助手来说，更稳妥的方式是先做分流：
\begin{itemize}
  \item \textbf{知识文本线}：原样保留事实与出处，优先服务 RAG。
  \item \textbf{任务样本线}：围绕目标任务构造输入输出对，服务 SFT。
  \item \textbf{语料适配线}：保留领域语言风格与结构分布，服务继续预训练。
\end{itemize}

这意味着“好数据”不是一个抽象标签，而是\textbf{与用途绑定的好}。适合进检索库的 PDF，不一定适合直接做 SFT；适合做继续预训练的海量规范文本，也不一定能直接作为高价值问答样本。

\subsection{RAG 语料和训练语料，不要默认共用同一池}
很多团队在数据工程里踩的一个大坑，是把“能喂给 RAG 的文本”和“能喂给训练的文本”当成同一种资源。表面上看，它们都来自企业文档；但在系统职责上，它们完全不是一回事。RAG 语料承担的是\textbf{实时事实来源}，强调版本新、出处清、可引用；训练语料承担的是\textbf{行为塑形或语言分布适配}，强调样本边界稳、信号一致、能被长期复用。

因此，同一份制度文档可以一边进入 RAG 库，一边被抽取少量高价值片段做 SFT 或继续预训练候选，但绝不能因为“它已经在知识库里”就默认“它也适合直接进训练池”。一份适合检索的长 PDF，可能充满页眉页脚、附录、版本差异和历史例外条款；如果不经切片和筛选直接拿去训练，模型学到的往往不是制度本身，而是制度文件的排版噪声、历史歧义和跨版本混杂表达。

更实用的做法是强制区分两个问题：
\begin{enumerate}
  \item 这份材料是否应该被检索到，供模型在回答时引用。
  \item 这份材料里是否有某些片段值得被写进训练样本，去塑造回答习惯或语言分布。
\end{enumerate}

只要这两个问题在流程上被分开，很多“为什么模型把旧制度写死在参数里”“为什么检索库一更新，训练集也跟着乱了”的问题就会自然减少。

\subsection{SFT 数据的三条硬标准}
SFT 数据不是把企业文本改写成问答对就够了。它至少要满足三条硬标准：
\begin{enumerate}
  \item \textbf{任务对齐}：样本必须直指目标任务，而不是“看起来像问答”。
  \item \textbf{答案质量}：输出要体现团队希望模型学习的风格、边界和证据使用方式。
  \item \textbf{格式稳定}：如果期望模型输出“结论-依据-建议动作”，训练样本就必须稳定体现这一规范。
\end{enumerate}

很多低质量 SFT 数据的问题，不是答案完全错误，而是\textbf{把模糊、含糊、风格混乱的表达也当作示范答案}。模型会忠实学习这种不稳定，而不是自动替团队补齐规范。

\subsection{拒答样本、边界样本和澄清样本，不是可选配菜}
如果企业知识助手未来要在线上面对真实用户，SFT 数据里就不能只有“有标准答案的问题”。教材里要把这一点讲得更绝对一些：\textbf{只教模型怎么答，不教模型什么时候不该答、什么时候该追问、什么时候该升级人工，等于把系统最贵的错误留到上线后再交学费。}

对工程读者来说，至少要把下面三类样本单独建池并显式覆盖：
\begin{itemize}
  \item \textbf{拒答样本}：材料不足、权限不足、问题越权、版本过期时，模型该如何明确说“不能判断”。
  \item \textbf{边界样本}：材料存在，但带有例外流程、跨部门责任或合规前提时，模型该如何给出保守答案，而不是擅自补全。
  \item \textbf{澄清样本}：用户问题本身含混，例如“这个能不能报”“这个流程过了吗”，模型该先问哪一个关键澄清问题。
\end{itemize}

这三类样本的价值，不在条数，而在它们决定了模型会不会形成正确的风控习惯。很多团队后期花大量时间补“安全规则”“兜底提示”，本质上是在用系统层补偿训练层的空缺。若第 6 章不把这件事写清楚，读者很容易误以为数据工程只负责堆“能答对”的样本，而不是同时建设“该保守时保守”的能力。

\subsection{领域语料来源优先级}
继续预训练不要求每条语料都有标准答案，但非常依赖语料整体质量。旧稿里最有价值的一条经验是：\textbf{优先选择知识密度高、语言质量稳、来源可控的文本}，而不是为了规模先抓一切能抓到的内容。

\begin{table}[H]
\centering
\small
\begin{tabularx}{\textwidth}{>{\raggedright\arraybackslash}p{3cm}>{\raggedright\arraybackslash}p{2.2cm}>{\raggedright\arraybackslash}p{2.2cm}>{\raggedright\arraybackslash}X}
\toprule
\textbf{语料来源} & \textbf{知识密度} & \textbf{可控性} & \textbf{建议} \\
\midrule
标准、制度、手册、教材、技术文档 & 高 & 高 & 继续预训练优先来源 \\
清洗后的高质量网页、官方博客、产品文档 & 中高 & 中 & 可作为补充语料 \\
论坛讨论、社交媒体、低门槛问答社区 & 波动大 & 低 & 谨慎使用，通常不作为主语料 \\
原始聊天记录、工单对话、会议纪要 & 中 & 中 & 更适合做筛选后样本，不宜整批无差别入池 \\
\bottomrule
\end{tabularx}
\caption{继续预训练语料来源的优先级}
\label{tab:cpt-source-priority}
\end{table}

\subsection{数据不是越多越好，而是越可控越好}
对工程项目而言，数据价值往往先取决于边界，再取决于规模。下面三件事比“先凑更多文本”更重要：
\begin{itemize}
  \item 文本是否来自可信、可授权的来源。
  \item 样本是否真的对应目标任务，而不是把噪声一并带入。
  \item 数据是否能在版本上被追踪，出了问题能回滚。
\end{itemize}

继续预训练尤其如此。大量格式混乱、时间过期、事实互相矛盾的语料，不是“先练练看”，而是会先把模型污染，再增加后续定位问题的成本。

\subsection{数据清洗最好拆成五道离线闸门}
很多团队把数据清洗理解成“去乱码”，但对训练项目来说，更合理的理解是：在样本入池前，先通过几道稳定的离线闸门，把明显会污染训练信号的内容挡在外面。这样做的好处是，训练时就不必一边烧卡、一边临时修数据。

一个对工程团队足够实用的最小流程，通常包括五道闸门：
\begin{enumerate}
  \item \textbf{来源与授权闸门}：确认文本来自可用、可追责、可回溯的来源。
  \item \textbf{格式标准化闸门}：统一编码、换行、标题层级、表格转写和特殊符号处理。
  \item \textbf{噪声移除闸门}：删除广告、导航、版权尾注、模板页眉页脚和明显脏片段。
  \item \textbf{重复控制闸门}：做文档级精确去重、段落级模糊去重和语义级近似去重。
  \item \textbf{质量与风险闸门}：过滤低信息密度、明显过期、相互冲突或含敏感风险的内容。
\end{enumerate}

\begin{table}[H]
\centering
\small
\begin{tabularx}{\textwidth}{>{\raggedright\arraybackslash}p{2.9cm}>{\raggedright\arraybackslash}p{3.2cm}>{\raggedright\arraybackslash}X}
\toprule
\textbf{闸门} & \textbf{主要动作} & \textbf{不过闸最容易出什么问题} \\
\midrule
来源与授权 & 标记来源、版本、生效时间、授权状态 & 后续无法回滚，甚至连“这批数据从哪来”都说不清 \\
格式标准化 & 统一编码、标题层级、表格与列表表示 & 同一类文本被切成不同形态，训练信号变散 \\
噪声移除 & 去广告、去导航、去模板噪声与脏页脚 & 模型学到无关格式、口水词和网页残留结构 \\
重复控制 & 精确去重、模糊去重、语义近重过滤 & 少数模板被过度放大，模型更像在背稿 \\
质量与风险 & 过滤低质、过期、冲突、脱敏不完整的文本 & 训练分数看似正常，但上线后边界和事实一起变脆 \\
\bottomrule
\end{tabularx}
\caption{数据清洗的目标不是“看起来干净”，而是让训练信号稳定可信}
\label{tab:data-cleaning-gates}
\end{table}

把这五道闸门放在离线阶段，还有一个经常被低估的收益：\textbf{每一道闸门都可以独立回放、独立复查、独立修补。} 这样当模型突然学坏时，团队能先问“是哪一道闸门放进来了问题样本”，而不是只能笼统怀疑“数据好像有点脏”。

\subsection{脱敏与权限隔离不是法务附录，而是训练准入条件}
企业数据工程里还有一类问题，旧稿里虽然没有系统收束，但在正式教材版里必须直接讲透：\textbf{脱敏、授权与权限隔离不是训练之后再补的合规流程，而是数据能不能入池的前置条件。} 一旦把带个人信息、敏感字段、越权材料或带工号映射关系的文本直接送进训练，后果往往不是某条样本脏了，而是模型把不该学的内容扩散成参数级风险。

这类风险最麻烦的地方在于，它不像乱码那样显眼。很多时候，一条工单、一段客服对话、一次审批记录在业务侧看起来都“很有价值”，但如果没有先做脱敏和权限判定，模型学到的就可能不是任务边界，而是某个真实人的手机号、内部地址、审批习惯甚至敏感决策口径。因此，\textbf{训练语料的“可用”不只取决于质量，也取决于它是否允许被长期内化。}

\begin{table}[H]
\centering
\small
\begin{tabularx}{\textwidth}{>{\raggedright\arraybackslash}p{3cm}>{\raggedright\arraybackslash}p{3.2cm}>{\raggedright\arraybackslash}X}
\toprule
\textbf{检查动作} & \textbf{至少要问什么} & \textbf{忽略后最常见的后果} \\
\midrule
个人信息脱敏 & 姓名、电话、邮箱、身份证件、地址是否已处理 & 模型记住不该复现的个体信息 \\
权限分层 & 文本是否只面向特定岗位、特定部门或审批链条 & 模型学会越权回答或泄露内部口径 \\
授权确认 & 来源是否允许进入训练与长期存储 & 项目后期无法合法复用或发布 \\
敏感标签登记 & 财务、合规、人事、法务等高风险语料是否单独标记 & 风险样本混入主池，后续难以定向审计 \\
\bottomrule
\end{tabularx}
\caption{企业训练数据的准入，不只要看质量，还要看能否被合法且安全地内化}
\label{tab:data-privacy-gate}
\end{table}

\subsection{去重、质量过滤与污染控制要做在训练前面}
旧稿在数据清洗这一块其实给了很多实操经验。教材版里最值得保留的一条原则是：\textbf{去重、质量过滤和评测污染控制，不是训练后的善后工作，而是训练前就该做的准入机制。}

对工程团队来说，至少有三类去重需要分清：
\begin{itemize}
  \item \textbf{精确去重}：同一份文档被重复抓取、重复导出或多个路径同时入库。
  \item \textbf{模糊去重}：模板化文档、改标题复用、只改少量字段的重复样本。
  \item \textbf{语义去重}：FAQ 改写版、同义制度说明、重复答案的变体样本。
\end{itemize}

除此之外，还要单独警惕\textbf{训练-评测污染}。如果后续回归集、高风险保留集、线上典型故障样本不加隔离地回流到继续预训练或 SFT 中，模型分数可能会快速变好，但团队得到的只是“看熟了题”，不是能力真的长出来了。

\begin{table}[H]
\centering
\small
\begin{tabularx}{\textwidth}{>{\raggedright\arraybackslash}p{3cm}>{\raggedright\arraybackslash}p{3.1cm}>{\raggedright\arraybackslash}X}
\toprule
\textbf{治理动作} & \textbf{解决什么问题} & \textbf{不做会怎样} \\
\midrule
精确去重 & 避免无意义重复训练 & 浪费算力，放大少数来源权重 \\
模糊/语义去重 & 避免模板记忆和重复偏置 & 模型更容易背格式、背措辞而不是学边界 \\
质量过滤 & 删掉脏文本、乱码、低信息密度语料 & loss 波动变大，错误模式更难定位 \\
污染隔离 & 隔离开发集、回归集、保留集 & 分数虚高，版本比较失真 \\
\bottomrule
\end{tabularx}
\caption{数据治理的目标不是“更干净一点”，而是让训练信号更可信}
\label{tab:data-governance-actions}
\end{table}

\subsection{为什么要做样本拼接}
在自回归训练中，很多原始文本样本都比较短。如果直接逐条训练，计算会浪费在 padding 和边界切换上。因此常见做法是把多个短样本按 token 长度拼接成更长序列，以提高训练效率。

但样本拼接并不等于“随便拼”。工程上最重要的一点，是意识到拼接带来的不仅是吞吐收益，也有\textbf{伪上下文风险}。如果把本不相关的问答、日志和制度片段硬拼到一起，模型会把它们误当成有意义的邻接关系，学到错误的语义共现。

\subsection{样本拼接的三种常见策略}
旧稿对样本拼接讲得比较散，这里把它压缩成三类最有工程价值的策略：
\begin{enumerate}
  \item \textbf{随机拼接}：实现最简单，适合先解决 padding 浪费问题，但跨样本干扰最大。
  \item \textbf{带分隔或分段 mask 的拼接}：通过边界标记或掩码，降低前后样本互相污染的风险。
  \item \textbf{语义聚类拼接}：把相似文档或相似任务拼在一起，适合希望长上下文更“有意义”的场景，但也更容易带来重复和泄漏。
\end{enumerate}

\begin{table}[H]
\centering
\small
\begin{tabularx}{\textwidth}{>{\raggedright\arraybackslash}p{3cm}>{\raggedright\arraybackslash}p{2.8cm}>{\raggedright\arraybackslash}p{2.8cm}>{\raggedright\arraybackslash}X}
\toprule
\textbf{拼接策略} & \textbf{主要收益} & \textbf{主要风险} & \textbf{适用判断} \\
\midrule
随机拼接 & 实现简单、节省 padding & 伪上下文最明显 & 先做基线时使用 \\
边界标记或分段 mask & 降低跨样本干扰 & 预处理复杂度更高 & 希望兼顾吞吐和干净边界时使用 \\
语义聚类拼接 & 局部长上下文更连贯 & 重复、泄漏、模板记忆风险更高 & 只有当相邻样本确实相关时使用 \\
\bottomrule
\end{tabularx}
\caption{样本拼接不是单一技巧，而是一组取舍}
\label{tab:packing-strategies}
\end{table}

\subsection{样本边界怎么隔离，比“是否拼接”更影响副作用}
旧稿里关于样本拼接还有一个很实用的提醒：\textbf{拼接策略真正决定副作用大小的，常常不是“拼不拼”，而是“边界怎么隔离”。} 很多团队把拼接理解成纯吞吐优化，于是默认所有 token 都可以互相看见，结果模型学到了一堆并不存在的上下文关系。

工程上至少要做三层判断：
\begin{itemize}
  \item \textbf{是否显式插入边界标记}：适合让模型知道“上一段已经结束”，避免把前一条工单结论延续到下一条问答。
  \item \textbf{是否限制跨样本注意力}：当样本彼此独立、格式又很强时，分段 attention mask 或近似的噪声屏蔽通常比单纯插分隔符更稳。
  \item \textbf{是否允许“有目的的相邻”}：如果确实在做语义聚类拼接，就要承认自己在给模型制造局部上下文，此时必须额外防重复、防泄漏，不能再把它当成中性的预处理。
\end{itemize}

可以把旧稿里提到的随机拼接、NoiseMask 和上下文式预训练理解成同一条轴上的三个点：最左边是“只求吞吐”，最右边是“希望相邻文本本身有学习意义”。越往右走，越不是简单的数据工程，而是在主动塑造模型会把什么当作同一段上下文。

\subsection{从原始材料到 SFT 样本，最好走四段流水线}
旧稿在 Self-Instruct 和 SFT 样本生成部分其实给了一个很重要的方法论：\textbf{不要把“生成样本”当成一次性动作，而要把它拆成四段流水线。} 这样做的好处是，每一段都能单独质检、单独回滚。

\begin{enumerate}
  \item \textbf{任务切片}：先把原始材料切成可回答的任务单元，例如制度问答、步骤解释、字段抽取、摘要改写，而不是一上来就整段喂给模型。
  \item \textbf{任务类型识别}：判断这条材料更像分类、抽取、问答还是流程建议。类型不清，后面的提示模板和输出格式就很难统一。
  \item \textbf{输入输出生成}：可以人工写、可以模型扩写、也可以先写答案再反推问题，但必须明确这一轮是在扩覆盖面，还是在固化风格。
  \item \textbf{后处理与抽检}：去重、脱敏、格式统一、风险场景抽检、与评测集隔离，最后再决定是否入库。
\end{enumerate}

以企业报销制度为例，同一份原始材料可以被切成很多种样本：一类是“发票缺字段时该怎么回复”的问答样本；一类是“把制度条款摘要成审批建议”的结构化输出样本；还有一类是“证据不足时必须升级人工确认”的保守回答样本。它们都来自同一份原文，但生成逻辑、验收标准和后续用途并不相同。

这也是为什么旧稿里反复强调“任务类型识别”不是多余步骤。若把本该做抽取的数据按开放问答生成，模型就会学得太散；若把本该学保守边界的样本按单一模板扩写，模型又会学得太硬。真正稳定的数据流水线，不是让模型写得越多越好，而是让每一条新增样本都知道自己在补哪一块能力缺口。

\subsection{表格、日志与流程文本要先正规化成可训练单元}
企业知识助手面对的原始材料，往往并不都是规整自然语言。它还可能包含字段表、审批流程、工单日志、错误码列表、混合中英术语、以及被复制过多次的半结构化模板。旧稿里这些内容分散在数据集构建、继续预训练和样本生成章节里；新版更适合给出一个统一判断：\textbf{异构材料如果不先正规化，就不该直接进入训练。}

这里的“正规化”不是把所有内容都改写成散文，而是把它们整理成模型能稳定理解、团队也能稳定回溯的训练单元。比如，表格要保住字段名与列对齐关系，日志要保住时间顺序与关键状态迁移，流程文本要保住步骤边界和分支条件。如果在这一步就把结构打碎，后面再多的训练也只是放大噪声。

\begin{table}[H]
\centering
\small
\begin{tabularx}{\textwidth}{>{\raggedright\arraybackslash}p{3cm}>{\raggedright\arraybackslash}p{3.2cm}>{\raggedright\arraybackslash}X}
\toprule
\textbf{原始材料} & \textbf{正规化重点} & \textbf{最怕丢什么} \\
\midrule
制度表格与字段清单 & 保住列名、字段含义、适用条件 & 列关系被打散后，模型只剩零碎名词 \\
工单或系统日志 & 保住时间顺序、状态变化、关键事件 & 前后状态断裂，模型学不到因果链 \\
审批流程或 SOP & 保住步骤边界、分支条件、角色职责 & 模型只记口号，看不出流程触发条件 \\
术语词典与缩写表 & 保住术语、释义、别名映射 & 缩写和解释脱钩，领域理解仍然混乱 \\
\bottomrule
\end{tabularx}
\caption{异构业务材料要先转成稳定训练单元，而不是直接整批灌入模型}
\label{tab:heterogeneous-normalization}
\end{table}

\subsection{SFT 样本格式先定协议，再谈规模}
很多教材会直接讨论“要准备多少条指令数据”，但工程落地时，先卡住团队的往往不是条数，而是\textbf{样本协议不统一}。同样叫一条 SFT 样本，有的团队存成单轮问答，有的存成多轮消息，有的把证据、标签和版本号写在元信息里，有的则全丢进自由文本。协议不统一，后面的拼接、抽检、回归和故障定位都会变乱。

对企业知识助手来说，SFT 样本至少应该把“要学什么行为”和“凭什么这样回答”分清。一个够用的最小协议通常包括下面几项：
\begin{table}[H]
\centering
\small
\begin{tabularx}{\textwidth}{>{\raggedright\arraybackslash}p{2.6cm}>{\raggedright\arraybackslash}p{3cm}>{\raggedright\arraybackslash}X}
\toprule
\textbf{字段} & \textbf{作用} & \textbf{最少要求} \\
\midrule
instruction / system & 说明任务边界与角色 & 不要写成泛泛口号，要能约束回答方式 \\
input / context & 给出问题、材料片段或待处理对象 & 明确这是原始用户问题、制度片段还是工单上下文 \\
output & 作为示范答案或结构化目标 & 体现团队希望模型学习的格式、证据使用和保守边界 \\
metadata & 支撑治理与回溯 & 至少保留来源、任务类型、版本号和风险标签 \\
\bottomrule
\end{tabularx}
\caption{SFT 样本先统一协议，后续的质检和回滚才有抓手}
\label{tab:sft-schema}
\end{table}

如果是多轮对话样本，还要额外回答两个问题：哪些历史轮次必须保留，哪些系统提示只该出现一次；证据片段是并入对话正文，还是单独作为上下文字段。只有这些约定先写清，团队后面才不会一边训练，一边为“同一种任务为什么有三种格式”返工。

这里还要特别划清边界：\textbf{本章讲的是 SFT 样本协议，不把偏好对、排序对和 PPO 轨迹混进来。} 偏好数据属于第 10 章的对齐链路；若现在就把不同训练目标的格式混在一个样本池里，读者会很难形成清晰的数据分工观。

\subsection{数据验收不是抽象质量，而是一份入池合同}
很多团队会说“这批数据质量还可以”，但真到训练前，谁也说不清“可以”到底是什么意思。教材化写法里更好的表述是：\textbf{每一批数据在进入主训练池前，都应该通过一份可执行的入池合同。} 这份合同不追求华丽，而追求能让团队在出事时回到具体条款，而不是停留在感觉层面。

对企业知识助手来说，一份够用的入池合同至少应回答下面六个问题：
\begin{enumerate}
  \item 这批数据的来源、授权、版本和时间边界是否清楚。
  \item 它是服务 RAG、SFT 还是继续预训练，职责是否单一。
  \item 它有没有覆盖拒答、澄清、升级人工等关键边界行为。
  \item 它与开发集、回归集、保留集是否完成隔离。
  \item 它的格式协议、清洗规则和处理脚本是否可追溯。
  \item 它进入后最可能改善哪类行为，最可能带来哪类副作用。
\end{enumerate}

最后一条尤其重要。成熟的数据工程不只问“这批数据好不好”，还问“\textbf{它准备改哪一类错误}”。如果这个问题答不上来，那么就算样本数量再多，也更像往系统里继续加素材，而不是加训练信号。

\subsection{什么时候需要继续预训练}
继续预训练适用于这样一类场景：基础模型对通用中文和英文都不错，但对企业领域文本的语言分布不够熟悉，例如内部术语、日志格式、工单字段、制度表达和产品名词大量出现。

如果问题主要是“风格不一致”或“任务格式不稳定”，优先考虑 SFT；如果问题主要是“模型连领域文本都读不顺”，继续预训练才更合适。还要注意，继续预训练不是知识更新的默认方案。凡是高频变化、版本敏感的事实，优先还是应该通过 RAG 接入。

\subsection{继续预训练的数据量与混合策略}
继续预训练最容易被问到的问题之一，是“到底需要多少数据”。更实用的回答不是报一个固定数字，而是看三件事：
\begin{itemize}
  \item 目标语料与基座模型原始分布差距有多大。
  \item 企业术语和结构化表达的密度有多高。
  \item 目标是让模型“读顺文本”，还是同时承担下游行为迁移。
\end{itemize}

如果领域数据量不大，完全只喂垂域语料，反而可能导致灾难性遗忘。因此工程上常见做法是把领域语料与通用语料混合训练。旧稿里反复强调的一条经验是：可以把领域语料与通用语料按大约 \(1:5\) 到 \(1:10\) 作为起始尝试，但这只是经验起点，不是固定法则。真正重要的是持续看回归评测：领域能力有没有提升，通用能力有没有明显塌陷。

\subsection{混比不是常数，而是一条阶段曲线}
继续预训练里另一个很容易被简化过头的问题，是把“领域语料和通用语料的比例”理解成一个固定常数。更成熟的工程理解是：\textbf{混比往往要跟阶段一起变化。} 起步阶段，团队更需要稳住基座分布，避免一上来就被垂域语料带偏；进入中段后，才逐步提高领域语料权重；到了末段，如果已经确认要衔接后续任务训练，还可能少量混入任务样本或高价值模板样本。

把混比看成阶段曲线，而不是一次拍板的数字，有两个直接收益。第一，团队可以更清楚地解释某一轮收益到底来自“更懂领域语言”还是“更像下游任务”；第二，当通用能力回退时，也更容易回滚到某一阶段的配比，而不是重新从头猜比例。

\begin{table}[H]
\centering
\small
\begin{tabularx}{\textwidth}{>{\raggedright\arraybackslash}p{2.8cm}>{\raggedright\arraybackslash}p{3.1cm}>{\raggedright\arraybackslash}X}
\toprule
\textbf{训练阶段} & \textbf{混比思路} & \textbf{主要目标} \\
\midrule
起步阶段 & 通用语料占更高比例，领域语料保守引入 & 先稳住基座行为，不让模型突然偏科 \\
主体阶段 & 逐步提高领域语料权重 & 让模型真正熟悉企业术语、格式和表达分布 \\
收束阶段 & 视目标少量混入任务样本或高价值模板 & 为后续 SFT 或任务适配做平滑衔接 \\
\bottomrule
\end{tabularx}
\caption{继续预训练的数据混比更像阶段调度，而不是一个永远不变的常数}
\label{tab:data-mixture-schedule}
\end{table}

\subsection{MIP：继续预训练末段少量混入任务样本}
旧稿里有一个容易被忽略、但对工程路径很有帮助的概念：\textbf{Multi-Task Instruction PreTraining}，简称 MIP。它的直觉并不复杂，就是在继续预训练的后段少量混入任务指令样本，让模型在适应领域语言分布的同时，不至于与后续任务行为训练完全脱节。

它适合的场景通常是：
\begin{itemize}
  \item 你已经确认需要继续预训练，而不是只做 SFT。
  \item 你手里有一小批质量不错的任务样本，但数量还不足以单独支撑完整 SFT。
  \item 你希望继续预训练后的模型，不至于在下游任务上“语料吃熟了，但行为还没接上”。
\end{itemize}

但 MIP 不是把两种训练目标随便混在一起。工程上更稳的做法是：\textbf{继续预训练仍以领域语料为主，任务样本只占很小比例，更多承担过渡和缓冲作用。} 如果让任务样本反客为主，就会把继续预训练重新拖回任务微调，反而模糊了本章最重要的分工边界。

\subsection{词表扩增通常不是领域适配的第一刀}
旧稿还专门提过“词表扩增必要性”。这在新版里也应该明确保留一个结论：\textbf{词表扩增更多影响分词形态与解码效率，通常不是多数领域适配问题的首要矛盾。}

只有在下面这类情况里，它才值得认真评估：
\begin{itemize}
  \item 内部缩略语、代码、字段名被切得极碎，显著影响 token 利用效率。
  \item 长表格字段、日志模式或混合语言术语反复被异常拆分。
  \item 你已经验证过：问题不只是训练数据不足，而确实与分词形态强相关。
\end{itemize}

否则，更常见、更高收益的动作通常仍然是：先提高语料质量、改善继续预训练覆盖、收紧任务样本边界，而不是急着动词表。

\subsection{合成数据不是灌水：Self-Instruct、回译与人工种子}
旧稿里关于 SFT 数据生成也有不少可吸收的内容。新版里更适合把它收束成一个工程判断：\textbf{合成数据的价值不在“凑条数”，而在低成本扩充覆盖面，同时维持边界和风格。}

最常见的三条路线分别是：
\begin{itemize}
  \item \textbf{Self-Instruct}：先给少量种子任务，再让模型扩写更多指令-回答对，适合快速扩展任务覆盖面。
  \item \textbf{回译或逆向构造}：先准备高质量答案、解释或制度片段，再反推问题，适合希望控制答案质量时使用。
  \item \textbf{人工种子 + 模型扩写 + 人工抽检}：成本和质量最均衡，通常最适合真正要上线的企业场景。
\end{itemize}

\begin{table}[H]
\centering
\small
\begin{tabularx}{\textwidth}{>{\raggedright\arraybackslash}p{2.8cm}>{\raggedright\arraybackslash}p{3cm}>{\raggedright\arraybackslash}X}
\toprule
\textbf{合成路线} & \textbf{主要优势} & \textbf{主要风险} \\
\midrule
Self-Instruct & 扩样快，容易快速铺开任务面 & 风格单一，容易复制模型原有偏见 \\
回译构造 & 更容易保证答案侧质量 & 问题分布可能不自然，像“考试题”而不像真实用户问题 \\
人工种子 + 自动扩写 & 质量和规模更平衡 & 需要明确抽检与过滤流程，否则仍会积累噪声 \\
\bottomrule
\end{tabularx}
\caption{合成数据的关键不是能不能生成，而是能不能被可靠过滤}
\label{tab:synthetic-data-routes}
\end{table}

因此，合成数据最好只承担两种任务：第一，补人工样本难以覆盖的边角场景；第二，放大已经被人工确认过的风格和结构。若希望它直接替代高价值人工样本，往往会把模型已有的幻觉、偷懒和模板偏置再复制一遍。

\subsection{主动学习不是“高级玩法”，而是把标注预算花在刀刃上}
旧稿里关于主动学习和失败样本发现的部分，也值得在教材版里保留下来。对工程团队来说，主动学习不必理解成复杂算法，先理解成一个简单原则就够了：\textbf{不要平均标所有样本，而要优先标那些最可能改变模型行为的样本。}

这类样本通常有三种来源：
\begin{itemize}
  \item \textbf{模型最犹豫的样本}：多次采样结果分歧大，或多个候选回答优劣很难分。
  \item \textbf{业务最敏感的样本}：权限、合规、财务、客服升级等高风险问题，即使数量少，也值得优先标。
  \item \textbf{历史上最容易出错的样本}：来自线上失败会话、人工驳回记录或评测回归里的顽固错误。
\end{itemize}

如果把这三类样本持续回流到 SFT 和评测流水线里，团队就不再只是“越标越多”，而是在逐步逼近最影响上线质量的边界区域。旧稿里提到的数据去重、差异性发现和不确定性采样，本质上讲的就是这件事：\textbf{用有限标注成本，优先覆盖模型最缺、最难、最贵的那部分行为。}

\subsection{合成样本也要做质量评分与长尾覆盖检查}
合成数据一旦进入主训练池，就不该再被当成“临时素材”，而要接受和人工样本类似的质量检查。更现实的做法，不是为每条样本打一个神秘总分，而是先建立几条能执行的准入问题：
\begin{itemize}
  \item 这条样本的答案是否真的正确，还是只是“像个答案”。
  \item 它是否和现有高频样本高度重复，进一步加重头部偏置。
  \item 它补的是长尾场景，还是只是在放大模型已经很熟的表达模式。
  \item 它的风格是否符合企业助手要求的保守边界、证据口吻和结构格式。
\end{itemize}

\begin{table}[H]
\centering
\small
\begin{tabularx}{\textwidth}{>{\raggedright\arraybackslash}p{2.5cm}>{\raggedright\arraybackslash}p{3.4cm}>{\raggedright\arraybackslash}X}
\toprule
\textbf{质检维度} & \textbf{最想回答什么问题} & \textbf{简单可执行的检查方式} \\
\midrule
准确性 & 答案是否站得住 & 人工抽检金标对比、规则校验、关键字段核对 \\
一致性 & 同类任务是否写得前后一致 & 同任务多样本交叉比对，检查格式和边界是否漂移 \\
多样性 & 是否只是重复头部样本 & 做去重与聚类，看相似样本是否过度堆积 \\
覆盖度 & 长尾场景有没有真的补上 & 对照任务清单检查冷门类别是否仍然空白 \\
边界性 & 模型是否学会该保守时保守 & 单独抽检拒答、升级人工、证据不足等高风险场景 \\
新鲜度 & 是否仍在反复放大旧问题分布 & 观察最近失败样本是否被及时纳入种子集 \\
\bottomrule
\end{tabularx}
\caption{合成样本的价值不在条数，而在是否补到了真实能力缺口}
\label{tab:synthetic-quality-gates}
\end{table}

这张表真正想传达的不是“再造一套复杂评审体系”，而是一个更朴素的工程纪律：\textbf{凡是自动生成得很快的东西，更要用明确的闸门限制它进入主训练池。} 否则，模型扩写带来的不是覆盖面，而是把已有偏差成倍复制。

\subsection{继续预训练中的训练信号}
虽然本章不展开完整训练工程，但有几个与数据强相关的训练信号必须在这里先建立直觉。

第一，\textbf{loss 突刺不一定意味着模型坏了}。当语料分布突然变化、脏数据混入、学习率过高，或从旧检查点接着训练时没有重新 warmup，都会出现不正常的 loss 上冲。第二，\textbf{重新 warmup 往往是继续预训练的必要动作}，因为此时模型面对的是一套新的语料分布，而不是完全延续原始预训练。第三，\textbf{batch size、学习率和样本拼接方式是联动的}，不能单看其中一个超参数。

\begin{table}[H]
\centering
\small
\begin{tabularx}{\textwidth}{>{\raggedright\arraybackslash}p{3cm}>{\raggedright\arraybackslash}p{3.5cm}>{\raggedright\arraybackslash}X}
\toprule
\textbf{观察到的现象} & \textbf{更可能的原因} & \textbf{第一动作} \\
\midrule
loss 突然上冲 & 语料分布切换太猛、脏数据混入、学习率过高 & 先抽查最近入池语料，再核对学习率与 warmup 设置 \\
领域开发集不涨 & 数据覆盖仍太散，没补到关键术语和结构模式 & 回到数据分流与样本切片，确认是不是还在喂泛泛文本 \\
通用能力明显回退 & 领域语料比例过高，或训练时长已经越界 & 增加通用混合语料，缩短训练并查看哨兵集趋势 \\
训练 loss 下降，但下游任务没改善 & 学到的是语料表面模式，不是可迁移行为 & 追加下游回归集，检查是否应转向 SFT 或 RAG \\
\bottomrule
\end{tabularx}
\caption{继续预训练里的很多异常，首先是数据和节奏问题，不只是模型问题}
\label{tab:cpt-signal-diagnosis}
\end{table}

\subsection{继续预训练也要有开发集和停机标准}
旧稿里不少实验分析都在提醒同一件事：\textbf{继续预训练不能只盯训练 loss，更不能靠“感觉读起来顺了一点”决定要不要继续烧卡。} 更稳的做法是同时准备三类验证信号：
\begin{itemize}
  \item \textbf{领域语料开发集}：看模型是否真的更会读内部术语、日志模式和制度表达。
  \item \textbf{下游回归小集}：看继续预训练后，摘要、问答、抽取等后续任务有没有被拖坏。
  \item \textbf{通用能力哨兵集}：看基础通用能力是否出现明显塌陷，避免“垂域变强，整体变脆”。
\end{itemize}

这三类信号缺一不可。只看领域开发集，容易忽视灾难性遗忘；只看下游任务，又容易把 SFT 收益误当成继续预训练收益；只看总 loss，则根本分不清模型是在学有用模式，还是在记脏数据和模板噪声。

工程上可以给继续预训练设一个很朴素的停机标准：当领域收益开始平台化、通用能力开始下滑、而后续 SFT 或 RAG 才是更合适的下一步时，就不该再靠继续预训练硬顶。教材里需要把这个判断留给读者，因为真正成熟的数据工程，不只是知道怎么开训，也知道什么时候该停、该换手段。

\subsection{数据版本化与质量闭环，本身就是训练资产}
如果一本教材只讲“如何构造样本”，却不讲“如何治理版本”，工程读者落地时仍然会卡住。因为真正导致项目不可复现的，往往不是模型参数，而是数据版本不清。

更稳的做法是给每一批训练数据都保留最少元信息：
\begin{itemize}
  \item 来源和授权状态；
  \item 清洗规则版本；
  \item 去重和过滤策略；
  \item 是否包含合成数据、占比多少；
  \item 与开发集、回归集、保留集的隔离关系。
\end{itemize}

这样做的价值有三层。第一，效果突然波动时，能回溯到底是样本变了还是训练变了。第二，某一批数据引入副作用时，可以回滚，而不是重新猜。第三，团队可以逐步建立“什么样的数据改动最值钱”的经验账本，而不是每次都从零开始试。

\subsection{数据漂移与新鲜度策略：不要等模型出错了才补}
数据治理还有一层经常被低估的内容：\textbf{新鲜度策略。} 企业知识助手服务的不是静态世界，制度会改版，产品字段会变，工单标签会扩展，用户提问方式也会随组织变化而漂移。若团队只是等线上错误积累到足够明显，再回头补数据，训练链路永远会落后于业务变化。

更成熟的做法，是给三条数据线分别设刷新节奏。知识文本线更关注版本变更和失效替换；任务样本线更关注新任务形态、新拒答场景和新格式要求；评测回流线更关注最近失败样本是否进入回归集。也就是说，\textbf{“更新数据”不是一个动作，而是三条线各自以不同频率、不同准则在更新。}

教材里把这点写出来，是为了帮助读者建立长期运营视角。数据工程不止是“把数据准备好让训练开起来”，更是“让系统在三个月后仍然知道哪些旧样本该下线，哪些新样本该补进来”。真正稳定的企业助手，靠的不是某一次训得多漂亮，而是这套新鲜度闭环能持续运转。

\section{主案例推进}
在企业知识助手中，我们把数据建设分成三条线并行推进。

\textbf{第一条线是知识文本线}：收集制度、产品文档、操作手册、FAQ 和工单模板。这些数据一部分会进入 RAG 系统，一部分会作为继续预训练候选语料。

\textbf{第二条线是任务样本线}：围绕真实使用场景构造 SFT 数据，例如：
\begin{itemize}
  \item 根据制度材料回答“是否允许”“如何申请”“谁负责”类问题。
  \item 根据工单上下文生成分类标签和下一步建议。
  \item 把长文档摘要成规定格式。
\end{itemize}

\textbf{第三条线是评测回流线}：把线上失败会话、人工抽检样本和关键制度问答整理成评测集，用于后续版本回归。

例如，一段制度原文可以先进入知识文本线，供后续 RAG 检索；再从中抽取一个高价值问答对，进入任务样本线，教模型学会“先给结论，再给依据”的回答形式；最后还可以从中挑出一批带版本号和标准答案的问题，进入评测回流线。也就是说，同一份原始材料可能服务于不同数据资产，但不能用同一种加工方式处理到底。

把这套流程写成工程步骤，会更清楚：
\begin{enumerate}
  \item \textbf{来源审计}：确认版权、授权、时效性和版本号。
  \item \textbf{文本清洗}：去广告、去模板噪声、统一编码、去重、过滤低质量片段。
  \item \textbf{资产分流}：按用途进入 RAG、继续预训练或 SFT 流水线。
  \item \textbf{版本登记}：记录数据来源、切分方式、处理脚本和模型版本。
  \item \textbf{回归验证}：训练后用独立评测集看收益与副作用。
\end{enumerate}

\subsection{给三条数据线分别设刷新节奏}
主案例如果只讲“怎么分三条线”，还不够像一本教材；更重要的是要告诉读者，这三条线应如何持续运转。对企业知识助手来说，可以先采用一套非常朴素但有效的节奏：
\begin{itemize}
  \item \textbf{知识文本线}：跟着制度发布和产品版本走，一旦文档改版，优先更新检索资产和引用版本标记。
  \item \textbf{任务样本线}：跟着高频失败类型走，每周或每两周补一轮“新边界、新拒答、新格式”的样本。
  \item \textbf{评测回流线}：跟着线上事故和人工抽检走，把最近最贵的错误优先固化成回归题。
\end{itemize}

这套节奏看起来不复杂，却能把“数据建设”从一次性工程改成持续运营。知识文本线保证模型说话时有最新依据，任务样本线保证模型学到新的回答习惯，评测回流线保证团队能看见版本是否真的变好。三条线只要有一条停了，系统都会逐渐失去弹性：要么知识过期，要么行为僵化，要么团队再也看不清自己到底在进步还是在回退。

下面给出一个适合工程团队使用的数据质检表：

\begin{longtable}{>{\raggedright\arraybackslash}p{3.2cm}>{\raggedright\arraybackslash}p{3.8cm}>{\raggedright\arraybackslash}p{6.2cm}}
\toprule
\textbf{检查项} & \textbf{适用数据} & \textbf{通过标准} \\
\midrule
来源可信与授权 & 全部数据 & 来源明确，可确认使用权限 \\
时效性 & 知识文本、评测集 & 能标出版本或生效时间 \\
任务对齐 & SFT 数据 & 样本直接对应目标任务，不是泛泛文本 \\
结构完整 & 表格、制度、工单 & 关键信息没有在切分中丢失 \\
版本可追踪 & 全部数据 & 能定位数据变更与模型版本对应关系 \\
训练/评测隔离 & SFT、评测集 & 评测样本不回流进训练主集 \\
\bottomrule
\end{longtable}

如果要把“原始材料去向判断”再做得更细，可以参考下面这张分流表：

\begin{table}[H]
\centering
\small
\begin{tabularx}{\textwidth}{>{\raggedright\arraybackslash}p{3cm}>{\raggedright\arraybackslash}p{3cm}>{\raggedright\arraybackslash}p{3.2cm}>{\raggedright\arraybackslash}X}
\toprule
\textbf{原始材料} & \textbf{默认去向} & \textbf{什么时候再转成 SFT} & \textbf{补充说明} \\
\midrule
最新制度 PDF & RAG / 评测线 & 需要固定回答格式时 & 动态知识优先保留为可检索文本 \\
产品手册与操作手册 & RAG / 继续预训练 & 有高频标准问答时 & 语言分布稳定，适合进语料池 \\
工单与客服对话 & SFT 候选 / 评测回流 & 经过清洗和脱敏后 & 不宜原样整批进继续预训练 \\
内部术语词典、字段说明 & 继续预训练 / RAG & 需要统一解释风格时 & 适合缓解术语读不顺问题 \\
\bottomrule
\end{tabularx}
\caption{企业原始材料到不同数据资产的分流}
\label{tab:asset-routing}
\end{table}

再看一个样本拼接示例。若有三条很短的制度问答样本，直接逐条训练会造成大量 padding；更合理的做法是按 token 长度把它们拼成一个较长序列，并在段落之间插入清晰边界标记。这样既能提高训练效率，也能降低模型误把前一个问答的风格延续到下一个样本中的风险。

如果团队已经确定要做继续预训练，还可以采用一种更贴近真实工程的两阶段路线：先用高质量领域文本做继续预训练，让模型适应企业语言分布；再用任务样本做 SFT，让模型学会回答边界与输出格式。某些场景下，也可以在继续预训练后段少量混入任务指令样本，作为连接二者的缓冲层，但前提是不能让任务样本反客为主，破坏继续预训练的语料目标。

\section{常见坑与工程取舍}
\begin{itemize}
  \item \textbf{把所有企业文本都扔进继续预训练}：如果语料噪声高、格式乱，会先污染模型再谈提升。
  \item \textbf{训练集和评测集来源不分离}：得到的“提升”很可能只是记忆而非泛化。
  \item \textbf{SFT 数据只追求数量}：低质量自动生成数据会把模型推向模板化和幻觉化。
  \item \textbf{忽略版本管理}：当回答质量突然变化时，如果无法定位是模型版本还是数据版本导致，就难以迭代。
\end{itemize}

再补四个旧稿里反复出现、但很容易被压掉的提醒：
\begin{itemize}
  \item \textbf{只喂垂域语料会放大遗忘风险}：领域能力可能上来，但通用能力和稳健性会塌。
  \item \textbf{语义聚类拼接不是越像越好}：相似文本堆得太密，容易让模型记住模板而不是学会泛化。
  \item \textbf{不要把扩词表当领域适配捷径}：除非确实证明分词切分极差，否则词表扩展未必值得。
  \item \textbf{继续预训练不看训练信号}：loss 突刺、学习率过大、重启后不 rewarmup，都会把好语料练坏。
\end{itemize}

\section{本章练习}
\begin{exercise}[样本去向判断]
一份最新发布的差旅制度 PDF，应该优先进入哪条数据线？为什么不应默认先做 SFT 样本？
\end{exercise}

\begin{solution}
应优先进入{\heiti 知识文本线}，因为它首先承担“被检索、被引用”的职责。只有当团队明确希望模型学习某种回答格式或高频问答行为时，才从中抽取高价值片段构造 SFT 样本。否则直接做 SFT 会把动态知识过早写进参数。
\end{solution}

\begin{exercise}[继续预训练触发条件]
什么症状更能说明企业知识助手需要继续预训练，而不是只做 SFT？
\end{exercise}

\begin{solution}
如果模型{\heiti 连领域文本本身都读不顺}，例如内部术语、日志格式、工单字段理解明显吃力，这更像领域语言分布不熟，需要继续预训练。若只是回答格式不稳或任务边界不清，则优先做 SFT。
\end{solution}

\section{延伸阅读}
\begin{itemize}
  \item \textbf{DeepLearning.AI / Hugging Face 的 LLM 数据课程资料}：适合理解数据清洗、切分与格式化的基本套路。
  \item \textbf{Together、MosaicML、Databricks 的继续预训练实践文档}：适合理解领域适配数据怎样进入训练流水线。
  \item \textbf{Cleanlab、Argilla 等数据质检工具资料}：适合理解数据错误如何被系统化发现。
  \item \textbf{Self-Instruct 与主动学习相关文章}：适合理解任务样本如何从原始材料中被逐步构造出来。
\end{itemize}

\section{小结}
本章回答了“模型改造靠什么喂出来”的问题。继续预训练解决领域语言分布问题，SFT 解决任务行为问题，而高质量的数据流水线是二者的共同前提。真正成熟的数据工程，不是把文本越堆越多，而是让每一份原始材料都知道自己该去哪里、以什么形态进入系统。

\section{下一章过渡}
当模型、提示词和任务数据逐渐成形后，企业知识助手仍然面临一个关键难题：很多答案并不适合写入参数，而必须实时从文档中获取。下一章将正式进入 RAG 基础，构建检索增强问答系统的最小闭环。
